市场的共同点是认可锂价回升与 H1 利润大幅改善；分歧在于 H2 能否维持、雅化是否拥有足够的成本/资源优势，以及民爆利润能否抵消锂盐的周期波动。我们的预期差不是对 H1 预告本身更悲观，而是对将预告直接外推为全年高盈利更保守。

\section{外部预测：H1 已验证一部分，H2 假设仍高度分化}

三份完整公开券商 PDF 对 2026E 的归母预测为 23.15--30.32 亿元、EPS 为 2.01--2.63 元。东吴在 7 月 7 日预告后给出 42 元目标价和 16 倍 2026E PE；国信和开源的公开材料没有可用目标价。开源在本轮相对其前值上调 2026E/2027E 归母预测 51.2\%/66.0\%，说明外部预期正在随价格环境快速上修，也提高了下半年验证门槛。

\begin{exhibitbox}[表 2-1：公开卖方预测与 AStock 基准的差异]
\begin{tabularx}{\exhibitboxwidth}{L{2.2cm}R{2.1cm}R{2.1cm}R{1.9cm}X}
\toprule
机构 / 日期 & 2026E 收入 & 2026E 归母 & EPS & 目标价、方法与使用边界 \\
\midrule
国信 / 4 月 30 日 & 161.64 亿元 & 23.15 亿元 & 2.01 元 & 优于大市；未披露可用目标价 \\
开源 / 4 月 29 日 & 168.18 亿元 & 25.69 亿元 & 2.23 元 & 买入；未披露可用目标价 \\
东吴 / 7 月 7 日 & 170.52 亿元 & 30.32 亿元 & 2.63 元 & 买入；42 元，2026E PE 16 倍 \\
\textbf{AStock 基准} & \textbf{模型输出} & \textbf{20.1 亿元} & \textbf{1.74 元} & \textbf{17.0 元；自有周期校准多锚估值} \\
\bottomrule
\end{tabularx}
\end{exhibitbox}
\sourcenote{国信、开源、东吴完整研究报告；AStock 自有模型。外部预测不是公司指引，也不是 AStock 模型输入的替代品。}

基准归母较国信 2026E 低约 13\%，较东吴低约 34\%。差异主要来自对 H2 的持续性、现金转换和未披露成本优势的处理，而非否认 H1 的改善。外部预测对上行情景有参考价值，但在正式 H1 三表前不应以其目标价作为独立估值锚。

\section{我们的不同判断：把资源和客户的好故事留在牛情景}

\begin{exhibitbox}[表 2-2：预期差与可证伪条件]
\begin{tabularx}{\exhibitboxwidth}{L{2.8cm}X X}
\toprule
问题 & 市场/乐观叙事 & AStock 判断与证伪证据 \\
\midrule
H2 持续性 & Q2 高利润可延续全年 & 基准 H2 显著低于 Q2 年化；H1 收入、毛利和现金共同验证后才上修 \\
资源优势 & 资源布局可显著降低成本 & 未披露实际供给/成本/自供比例，基准零信用；披露可复算吨成本后再提高概率 \\
客户与订单 & 全球客户资格带来销量放量 & 无客户份额、量价和剩余订单，基准只认可渠道，不计增量收入 \\
民爆防线 & 民爆可完全平滑锂价波动 & 历史利润池可托底，但项目订单和回款未披露；现金质量决定其折现价值 \\
\bottomrule
\end{tabularx}
\end{exhibitbox}

\section{最强反方：若 H1 证明“利润、现金、成本”三者同步改善，当前估值可能过低}

反方并不依赖故事：若正式 H1 显示锂盐收入/毛利与销量、价格、成本一致，经营现金流显著修复，且民爆业务没有以更高营运资本换取利润，则市场对现金质量的折价会失去基础。届时东吴等更高盈利假设的概率应提高，17 元的基准目标会被重新校准。反方的强度恰恰要求我们现在不追逐未经验证的利润，而不是忽略它。
